\section{Resumen}
En las grandes ciudades, la disponibilidad de espacios de estacionamiento representa una problemática creciente debido al aumento del parque automotor y la limitada infraestructura urbana, generando congestión vehicular y dificultades para los conductores.

En este contexto, el presente proyecto propone el desarrollo de \textit{ParkLink}, una plataforma digital orientada a facilitar la búsqueda, publicación y reserva de espacios de estacionamiento, dirigida tanto a conductores como a propietarios de cocheras privadas y playas de estacionamiento.

La plataforma permitirá visualizar espacios disponibles mediante geolocalización, realizar reservas y gestionar publicaciones desde una única aplicación, incorporando autenticación de usuarios y procesamiento de pagos digitales.

En este trabajo desarrollaremos un Producto Mínimo Viable (MVP) que abarca los aspectos funcionales, técnicos y económicos del sistema, postergando integraciones avanzadas como sensores IoT o predicción con IA para fases posteriores.

\vspace{0.5cm}
\textbf{Palabras clave:} estacionamiento urbano, reserva de cocheras, geolocalización, plataforma digital, movilidad urbana.